% ABOUTME: Discussion and limitations section of the IEEE report.
% ABOUTME: Interprets the cross validation, test and explainability results and states the limits of the study.
\section{Discussion and Limitations}
\label{sec:discussion}

\subsection{Features and learners}
Both the feature set and the model family matter. Against the true test
labels, engineered features lower the RMSE by 2.22 cycles for XGBoost,
2.54 for Random Forest and 4.30 for Elastic Net, while replacing Elastic
Net with XGBoost lowers it by 4.99 on raw features and 2.91 on engineered
features; cross validation shows the same pattern. The paired bootstrap
interval for XGBoost on engineered minus raw features ($-4.4$ to $-0.2$
cycles) excludes zero, while the interval for Random Forest minus XGBoost
($-1.3$ to $+1.1$) does not, so once the features are fixed the choice
between the two tree ensembles makes no measurable difference on 100 test
engines. Our margin over the published Random Forest~\cite{zhang2017modbne}
comes mainly from the features: against labels clipped at 125 our Random
Forest scores 16.58 on raw sensors and 14.02 on engineered features,
against 17.91 published. The ablation of
Section~\ref{sec:features} places most of the feature gain on the rolling
slopes, which give the rate of degradation rather than its level.

\subsection{Accurate near failure, coarse far from it}
Both ensembles are accurate where maintenance is decided. On the 33 test
engines with true RUL of at most 50 cycles, the late-stage RMSE is 7.34
for Random Forest and 7.03 for XGBoost, against 17.88 and 18.11 on the 67
early-stage engines. Part of the early-stage error is built into the
target: eleven test engines have a true RUL above the training cap, and
no model on engineered features predicts above 130, so these engines are underestimated by
construction. This error only brings maintenance forward. The costlier
error appears between 51 and 75 cycles, where all three models predict too
late (median error $+9$ to $+14$ cycles), around the point at which
degradation first becomes visible in the sensors. Late predictions are the
error that the score $S$ penalises most.

\subsection{Model selection and the test set}
Cross validation selected XGBoost (RMSE 13.99 against 14.42 for Random
Forest), while Random Forest is marginally better on the test set (RMSE
15.23 against 15.36, score 344 against 389). We kept the selection made
before the test set was opened. The bootstrap interval above shows that the
test ranking of the two ensembles is within sampling noise, and most of the
score gap comes from a single engine that XGBoost predicts about 9 cycles
earlier than Random Forest (Section~\ref{sec:results}). Cross validation and test numbers are not
directly comparable: cross validation scores every cycle against the
capped target, whereas the test set scores one cycle per engine against
the true RUL. The cross validation scores of the selected configurations
are also optimistic, because the same folds chose and scored them.

\subsection{What the models rely on}
Held-out permutation importance shows that both ensembles depend mainly on
two inputs: the health index, whose shuffling raises RMSE by about 11
cycles for XGBoost and 7.4 for Random Forest, and the cycle number (4.84
and 5.22). No individual sensor group adds more than 1.62. The sensor
information therefore reaches the models largely through the health index,
a single learned combination of all 14 sensors. Reliance on the cycle
number means the models partly predict from age. This is useful on FD001,
where every engine starts in normal condition and training lifetimes span
128 to 362 cycles, but it ties the models to the lifetime distribution of
this fleet. Impurity and permutation importance also disagree for
correlated inputs: \texttt{s\_4} receives 26\% of Random Forest's impurity
reduction but costs only 1.62 cycles when shuffled, so impurity rankings
alone would overstate its role.

\subsection{Gap to deep sequence models}
Four published deep models are more accurate than our Random Forest and
XGBoost by 0.8 to 1.5 cycles of RMSE against clipped labels, and on the
score $S$ six published models do better than our Random Forest and eight
better than the selected XGBoost. These models learn
features from windows of the raw sensor sequence. Our features look back
at most 20 cycles and use window lengths chosen by hand, so they cannot
represent how long a sensor has been drifting or when its drift began to
accelerate. The remaining gap is consistent with this limit, although our
experiments do not isolate it.

\subsection{Limitations}
All results are on FD001, which has a single operating condition and a
single fault mode. The regime normalisation of Section~\ref{sec:features}
reduces to a z-score here and has not been tested on the multi-condition
subsets. The test set has 100 engines with one prediction each, so the
test metrics carry wide sampling intervals, as the bootstrap shows; the
late stage contains only 33 engines. The cap of 125 cycles bounds what the
models can predict, so engines with more remaining life are always
underestimated. Published FD001 results differ in protocol details that
are not always reported, such as whether test labels are clipped and the
window length used, so the comparison with them is approximate. Finally,
C-MAPSS is simulated, and conclusions about which sensors matter apply to
the simulation rather than to physical engines.
